









\section{An Observation on Loop Count and Reasoning-Step Length} 
\label{sec:appen:mechanistic_verification}



To generate hypotheses about why certain loop counts outperform others
(e.g., 9 vs.\ 11), we conduct an exploratory analysis of GSM8K responses
from Section~\ref{sec:expert}. We stress that the evidence below is
\emph{correlational}. We use MiniMax-M2.5 \citep{chen2026minimax} to
annotate responses following \citet{bogdan2025thought} and measure average
tokens per sentence within each semantic tag. The resulting distributions
(Figure~\ref{fig:avg_tokens_per_tag}) suggest a coupling between MTP
lookahead depth and the natural token length of reasoning steps, which we
term the \emph{horizon alignment} conjecture.


\begin{figure}[t]
    \centering
    \includegraphics[
        width=\linewidth,
        trim=0cm 9pt 0cm 12pt, %
        clip
    ]{figures/cross_model_tokens_per_tag.pdf}
    \vspace{-20pt}
    \caption{\textbf{Average number of tokens per annotated tag}. The reasoning traces are annotated following \citet{bogdan2025thought}.
    }
    \vspace{-10pt}
    \label{fig:avg_tokens_per_tag}
\end{figure}


\paragraph{Canonical sentence lengths define the alignment target.}
Across loop configurations, the model's reasoning vocabulary settles into a
narrow band of characteristic lengths: \texttt{active\_computation} (AC)
steps average ${\sim}14.5$--$15.5$ tokens, while
\texttt{problem\_setup} (PS) and \texttt{fact\_retrieval} (FR) steps cluster
at ${\sim}11$--$13$ tokens, and \texttt{plan\_generation} (PG) steps average ${\sim}14$--$16$ tokens. These values vary little with loop depth,
marking them as intrinsic structural units of the model's reasoning rather
than artifacts of a particular decoding configuration.

\paragraph{A possible account of the \texorpdfstring{$\text{Loop}15$}{loop 15} peak.} A 15-token lookahead is the smallest window that fully contains a
canonical AC step from opening clause to final result. Because the model can
preview an entire computational thought within a single horizon, it never
commits to an intermediate token without visibility into where the
calculation lands, explaining the accuracy peak at $\text{Loop}15$
($19.03\%$).

\begin{figure*}[t]
    \centering
    \includegraphics[
        width=\linewidth,
        trim=0cm 7pt 0cm 6pt, %
        clip
    ]{figures/grad_norm_and_spectral_norm.pdf}
    \vspace{-7pt}
    \caption{\textbf{Effect of WD on a looped model.} 
    (Left) Gradients during training for $T{=}9$ model
    (Right) Per-layer spectral norms of the trained model. 
    }
    \label{fig:spectral_norm}
\end{figure*}

\paragraph{An unexplained dip at \texorpdfstring{$\text{Loop}11$}{loop 11} and \texorpdfstring{$\text{Loop}13$}{loop 13}.} 
When the lookahead window neither spans a full execution step nor coincides with a stable
planning boundary, performance degrades sharply. At $\text{Loop}11$, the
horizon falls just short of the ${\sim}14.5$ token AC baseline; unable to preview the
outcome of a computation already begun, the model enters a recovery mode
visible as bloated \texttt{uncertainty\_management} (UM) sentences averaging
${\sim}28$ tokens---verbose, structurally unanchored attempts to re-derive
context the window could not supply. At $\text{Loop}13$, the horizon is wide
enough to \emph{enter} an AC step but too narrow to \emph{exit} it cleanly,
trapping generation in mid-computation limbo (${\sim}35$ UM tokens). Crucially, misalignment does
not suppress generation but redirects it into unproductive verbosity.


We stress the limits of this account. It offers a plausible reading of the $\text{Loop}15$ peak and the $\text{Loop}11$ / $\text{Loop}13$ dip, but it does not explain the strong results at $\text{Loop}7$ (18.12\%) and $\text{Loop}9$ (18.73\%), whose lookahead horizons are further from the canonical AC length than those of the configurations that underperform. A mechanism operating solely through horizon–step alignment would predict the opposite ordering for these points.






























